% =====================================================================
%  decision_layer.tex  --  PHASE III, Tasks 6 and 7
%
%  The optimisation-based shutdown decision layer: problem statement,
%  the reduction that makes it exactly solvable, the policy family,
%  and the empirical results including the economic sensitivity.
%
%  Intended placement: Section 5 (\label{sec:decision}), after the
%  forecasting section and before the sustainability analysis.
%
%  Drop-in usage: % =====================================================================
%  decision_layer.tex  --  PHASE III, Tasks 6 and 7
%
%  The optimisation-based shutdown decision layer: problem statement,
%  the reduction that makes it exactly solvable, the policy family,
%  and the empirical results including the economic sensitivity.
%
%  Intended placement: Section 5 (\label{sec:decision}), after the
%  forecasting section and before the sustainability analysis.
%
%  Drop-in usage: % =====================================================================
%  decision_layer.tex  --  PHASE III, Tasks 6 and 7
%
%  The optimisation-based shutdown decision layer: problem statement,
%  the reduction that makes it exactly solvable, the policy family,
%  and the empirical results including the economic sensitivity.
%
%  Intended placement: Section 5 (\label{sec:decision}), after the
%  forecasting section and before the sustainability analysis.
%
%  Drop-in usage: % =====================================================================
%  decision_layer.tex  --  PHASE III, Tasks 6 and 7
%
%  The optimisation-based shutdown decision layer: problem statement,
%  the reduction that makes it exactly solvable, the policy family,
%  and the empirical results including the economic sensitivity.
%
%  Intended placement: Section 5 (\label{sec:decision}), after the
%  forecasting section and before the sustainability analysis.
%
%  Drop-in usage: \input{decision_layer}
%  Bibliography: paper/references.bib
%  Numbers below are generated by  python run_phase3.py  and stored in
%  outputs/phase3/task6/<machine>/task6_results.json
%             and outputs/phase3/task7/<machine>/task7_results.json
% =====================================================================

\section{Optimisation-Based Shutdown Decision Support}
\label{sec:decision}

\subsection{Problem statement}
\label{sec:decision:problem}

Let $x(t)\in\{0,1\}$ indicate that the machine is energised at time $t$.
Over a horizon the plant incurs standby energy cost while idle and
energised, a restart cost each time the machine is re-energised, and a
production-delay cost because a cold machine reaches production more
slowly than a warm one. The shutdown policy is the solution of

\begin{equation}
\label{eq:objective}
\min_{x(\cdot)}\;
C_{\text{total}}
= c_{\text{sb}}\,h_{\text{standby}}
+ c_{\text{re}}\,N_{\text{restarts}}
+ c_{\text{dl}}\,h_{\text{delay}}
\end{equation}
subject to
\begin{align}
h_{\text{off}}^{(i)} &\ge h_{\min}
&&\text{(motor thermal cool-down)} \label{eq:c1}\\
N_{\text{restarts}}^{(d)} &\le N_{\max}
&&\text{(mechanical wear, per day $d$)} \label{eq:c2}\\
x(t) &\in \{0,1\}. \label{eq:c3}
\end{align}

The coefficients are
$c_{\text{sb}} = \pi\,P_{\text{sb}}$, with $\pi$ the electricity tariff
and $P_{\text{sb}}$ the mean standby power taken from the inferred
STANDBY component of Section~\ref{sec:method}; and
$c_{\text{re}} = \pi E_{r} + \ell + c_{\text{dl}} T_{d}$, with $E_{r}$
the restart energy, $\ell$ any non-energy per-cycle cost, and $T_{d}$
the restart delay. Section~\ref{sec:decision:coeff} measures $P_{\text{sb}}$,
$E_r$ and $T_d$ from the record itself.

\paragraph{The break-even duration is derived, not supplied.}
Setting the saving from de-energising an idle interval of length $D$ to
zero gives
\begin{equation}
\label{eq:breakeven}
D^{\star} \;=\; \frac{c_{\text{re}}}{c_{\text{sb}}}
\;=\; \frac{\pi E_{r} + \ell + c_{\text{dl}} T_{d}}{\pi\,P_{\text{sb}}} .
\end{equation}
This replaces the fixed threshold used in the conventional formulation,
in which $D^{\star}$ is an operator-supplied constant. Equation
\eqref{eq:breakeven} also carries a consequence that is easy to miss:
\emph{when the restart cost is purely electrical} ($\ell = 0$,
$c_{\text{dl}} = 0$) \emph{the tariff cancels}, leaving
$D^{\star} = E_{r}/P_{\text{sb}}$, and the break-even duration is
independent of electricity price. Price enters only through the
non-energy terms. We verify this numerically in
Section~\ref{sec:decision:sensitivity} and it governs how the
sensitivity analysis must be presented.

\subsection{Reduction to one decision per idle episode}
\label{sec:decision:reduction}

Although $x(t)$ is defined at every sample, $C_{\text{total}}$ changes
value only at the boundaries of an \emph{idle episode} --- a maximal
interval of non-productive operation bounded by production. Within one
episode the standby cost is linear and non-decreasing in the time spent
energised, so the saving from de-energising after $\tau$ seconds of an
episode of realised length $D$ is
\begin{equation}
\label{eq:epsave}
S(D,\tau) =
\begin{cases}
c_{\text{sb}}\,(D-\tau) - c_{\text{re}}, & \tau < D,\\
0, & \text{otherwise,}
\end{cases}
\end{equation}
which is decreasing in $\tau$. Any episode that is de-energised at all
is therefore optimally de-energised at its onset, and the problem
reduces to choosing one scalar $\tau_i$ per episode.

\begin{proposition}[Exact offline optimum]
\label{prop:oracle}
With realised durations known, the optimum of
\eqref{eq:objective}--\eqref{eq:c3} is obtained by selecting, within
each day, the $N_{\max}$ episodes of largest positive $S(D_i,0)$ among
those with $D_i \ge h_{\min}$.
\end{proposition}

\noindent
The objective is modular over episodes and the only coupling is the
per-day cardinality constraint \eqref{eq:c2}, so greedy selection is
optimal. Consequently the oracle bound reported below is the exact
solution of the stated problem rather than a heuristic for it, and no
dynamic programme over the $5.5\times10^{6}$ samples is required.

\subsection{Measuring the restart coefficients}
\label{sec:decision:coeff}

Rather than assume $E_r$ and $T_d$, we estimate them from the record.
Every resumption of production is classified by the state the machine
resumed \emph{from} --- cold (OFF) or warm (STANDBY) --- and the cost
attributable to having de-energised is the difference between the two.
Charging the entire cold-start ramp to the shutdown decision would bill
it for a ramp the machine would have paid in any case.

On pelletizer-I the machine reaches 90\% of production power in a median
$574$\,s from cold ($n=126$) against $308$\,s from warm ($n=682$), giving
\begin{equation}
T_{d} = 266\ \text{s}\qquad (\text{95\% CI } [140,\,351]),
\end{equation}
by percentile bootstrap. Energy integrated over a common $1800$\,s
window from the moment production resumes gives
$-6.21$\,kWh (95\% CI $[-7.99,\,-4.64]$; $n=109$ cold, $2526$ warm), and
a milestone-based estimator gives $-1.78$\,kWh.

Both energy estimates are negative and significantly so. This is not an
artefact: a cold start reaches production more slowly, so within any
fixed window it has produced less and therefore drawn less. The energy
difference and the delay are two views of one physical fact, and summing
a negative energy term with a positive delay term would double-count it.
We therefore report that \emph{this machine has no separately
identifiable electrical restart penalty --- its restart cost is a
production-delay cost} --- set $E_r = 0$ in the baseline, and sweep $E_r$
over the range reported in the literature in
Section~\ref{sec:decision:sensitivity}, which also covers machines for
which the penalty is real.

\subsection{Policies}
\label{sec:decision:policies}

Online, $D$ is unknown when the decision must be made. Paying a small
cost repeatedly or a large one-off cost that ends the payments is
exactly the ski-rental (spin-block) structure, and the classical
guarantee applies: de-energising once accumulated standby cost equals
the restart cost is $2$-competitive against the offline optimum for
\emph{any} duration distribution~\cite{karlin1994competitive,karlin1988competitive}.
The same structure underlies online device power-down~\cite{irani2003online}.
This --- not ``never shut down'' --- is the baseline a forecast-driven
policy must beat.

We compare: \textsc{observed}, the plant's own behaviour read off the
inferred labels; \textsc{never} and \textsc{immediate}, the two trivial
policies; \textsc{static}, the conventional rule that forecasts the
duration once at onset and compares it with $D^{\star}$;
\textsc{ski-rental}, the forecast-free elapsed-time rule; the proposed
\textsc{forecast-opt}; and \textsc{oracle} from
Proposition~\ref{prop:oracle}.

\paragraph{Proposed policy.}
Every $60$\,s while the machine remains idle, the rule re-decides using
a forecast of the \emph{remaining} idle time $R$:
\begin{equation}
\label{eq:rule}
\text{de-energise} \iff
\underbrace{c_{\text{sb}}\,\mathbb{E}[R \mid \mathcal{F}_\tau] > c_{\text{re}}}_{\text{economics}}
\;\wedge\;
\underbrace{\Pr(R \ge h_{\min} \mid \mathcal{F}_\tau) \ge 1-\alpha}_{\text{feasibility}} ,
\end{equation}
where $\mathcal{F}_\tau$ is the information available at elapsed time
$\tau$: elapsed idle time, the \emph{current} position in the factory
calendar, and the load history of the production run that just ended.

Two aspects of \eqref{eq:rule} are deliberate. First, it is
\emph{sequential} rather than a one-shot classification at onset --- the
moment when the least information is available --- so each additional
idle minute is used as evidence. Second, the two conditions are not
redundant. Because $S$ is affine in $R$ by \eqref{eq:epsave}, the
unconstrained risk-neutral rule depends on the predictive distribution
only through its \emph{mean}, and a mean regression suffices. Constraint
\eqref{eq:c1} breaks that affinity: a shutdown whose idle period proves
shorter than $h_{\min}$ is infeasible rather than merely unprofitable,
and no expectation over $R$ expresses this. Feasibility is therefore
imposed as a chance constraint~\cite{charnes1959chance} at a stated
confidence $1-\alpha = 0.90$, using a separate classifier for
$\Pr(R\ge h_{\min})$ (AUC $0.939$, Brier $0.102$). Setting $\alpha=1$
recovers the pure risk-neutral rule and isolates the constraint's
contribution.

\paragraph{The estimand matters more than the model.}
Idle durations span four orders of magnitude, so fitting
$\log(1+R)$ is the natural choice --- and it is the wrong one here.
Back-transforming a log-scale fit estimates the conditional
\emph{median}, and Duan's smearing correction~\cite{duan1983smearing}
repairs it with a single global constant that cannot absorb a
covariate-dependent bias. At episode onset the conditional median is
$7$\,s while the conditional mean exceeds $2\,000$\,s. Since
\eqref{eq:rule} consumes the mean, the estimand must be elicited by the
loss used to fit it~\cite{gneiting2011making}; we therefore fit the raw
scale under squared error and retain the log-scale model as an ablation.

\subsection{Experimental protocol}
\label{sec:decision:protocol}

Spike rows are retained: removing $19.5\%$ of a $1$\,Hz series destroys
the correspondence between row counts and seconds on which every episode
duration depends. Episodes touching one of the record's $31$
acquisition gaps are excluded ($51$ episodes, $187.9$\,h), because their
true extent is unknown and counting them would manufacture savings out
of missing data. This leaves $3\,224$ usable episodes spanning
$474.9$\,h of idle time. The split is \emph{chronological} ($70/30$):
unlike the unsupervised comparison of Section~\ref{sec:results}, a
decision policy is deployed forward in time, and a shuffled split would
let it learn the future duty cycle. The per-day restart cap is enforced
at evaluation time, identically for every policy. All annualisation uses
the $63.4$ days the record actually covers, not its $154$-day span.

\paragraph{The status quo is the baseline.}
Of the $474.9$ usable idle hours the plant already de-energises for
$388.1$, leaving $86.8$\,h energised. Savings quoted against a
``never shut down'' hypothetical would therefore be savings the plant has
already banked; every figure below is quoted against \textsc{observed}.

\paragraph{Scenarios.}
Every coefficient is measurable except the valuation of lost production,
and $D^{\star}$ is linear in it, so asserting a value would silently
determine the conclusion. Scenarios are instead indexed by where
$D^{\star}$ falls within the observed idle-duration distribution, and the
delay valuation each implies is reported as an \emph{output}
(Table~\ref{tab:decision:scenarios}).

\begin{table}[t]
\centering
\caption{Economic scenarios, indexed by the derived break-even point.
The implied production-delay valuation is an output of
Eq.~\eqref{eq:breakeven}, not an assumption.}
\label{tab:decision:scenarios}
\begin{tabular}{lrrrr}
\hline
Scenario & $D^{\star}$ & Implied $c_{\text{dl}}$ & $c_{\text{re}}$ & Episodes above \\
         & (min) & (USD/h) & (USD) & \\
\hline
S1 low       &  10 &  1.01 & 0.075 & 228 \\
S2 moderate  &  60 &  6.09 & 0.450 &  53 \\
S3 high      & 240 & 24.34 & 1.799 &  43 \\
\hline
\end{tabular}
\end{table}

\subsection{Results}
\label{sec:decision:results}

Table~\ref{tab:decision:policies} reports annual saving against the
status quo and the share of the oracle head-room captured, on $968$
held-out episodes spanning $78.3$ days.

\begin{table}[t]
\centering
\caption{Policy comparison on held-out episodes: annual saving versus
the plant's own behaviour (USD/yr) and share of the attainable head-room
(\%). The three middle rows are ablations of the proposed policy.}
\label{tab:decision:policies}
\begin{tabular}{lrrrrrr}
\hline
& \multicolumn{2}{c}{S1 (10\,min)} & \multicolumn{2}{c}{S2 (60\,min)} & \multicolumn{2}{c}{S3 (240\,min)}\\
Policy & USD/yr & \% & USD/yr & \% & USD/yr & \% \\
\hline
Never shut down            & $-215$ & $-530$ & $-133$ & $-190$ & $+163$ & $65$ \\
Observed (status quo)      & $0$ & $0$ & $0$ & $0$ & $0$ & $0$ \\
Immediate shutdown         & $-237$ & $-584$ & $-316$ & $-451$ & $-599$ & $-241$ \\
Static break-even          & $-98$ & $-242$ & $-46$ & $-65$ & $+175$ & $70$ \\
Ski-rental (no forecast)   & $-23$ & $-56$ & $+24$ & $34$ & $+106$ & $43$ \\
\quad ablation: median estimand      & $-23$ & $-56$ & $-14$ & $-20$ & $+180$ & $72$ \\
\quad ablation: no chance constraint & $-73$ & $-180$ & $+6$ & $8$ & $+173$ & $69$ \\
\textbf{Forecast optimisation}       & $\mathbf{-3}$ & $\mathbf{-8}$ & $\mathbf{+39}$ & $\mathbf{55}$ & $\mathbf{+190}$ & $\mathbf{76}$ \\
Oracle (offline optimum)   & $+41$ & $100$ & $+70$ & $100$ & $+249$ & $100$ \\
\hline
\end{tabular}
\end{table}

\paragraph{These figures are conditional on the labelling.}
Table~\ref{tab:decision:policies} is computed on the state labels of
Section~\ref{sec:states}, in which the Viterbi path is left
unconstrained. Re-deriving the same table from labels carrying the
minimum-dwell constraint of Section~\ref{sec:flicker}---the same
machine, the same policy, the same coefficients, and a test window cut
at the same instant---changes the S2 column materially: the attainable
head-room falls from $71$\,USD/yr to $38$\,USD/yr and the proposed
policy from $+41$\,USD/yr ($57\%$ of the oracle) to
$+2.9 \pm 0.6$\,USD/yr ($7.8 \pm 1.6\%$; mean and standard deviation
over ten random seeds from Phase~V Task~12C). The physical quantities
are unmoved (idle $474.9$ vs $474.1$\,h; energised idle $86.8$ vs
$87.8$\,h; standby power $3\,747$ vs $4\,208$\,W); what moves is the
\emph{partition of idle time into episodes}, which is the unit the
decision layer bids on: enforcing a minimum dwell merges $3\,224$
episodes into $1\,312$ and raises the median episode from $12$ to
$90$\,s. Sub-ten-second ``opportunities'' created by label flicker are
not shutdown opportunities, so the smaller figure is the defensible
one, and both are reported here rather than the more favourable one
alone. The seed-to-seed standard deviation ($0.6$\,USD/yr) is small
relative to the bootstrap sampling interval ($\pm 50$\,USD/yr over
18 factory days; Phase~V Task~11B), confirming that model-fitting
variance does not explain the uncertainty.

The proposed policy captures $55$--$76\%$ of the attainable head-room in
S2 and S3 and beats the forecast-free ski-rental rule in all three
scenarios. Both ablations are material: replacing the mean estimand with
the log-scale median costs $53$\,USD/yr in S2, and removing the chance
constraint costs $33$\,USD/yr in S2 while raising minimum-off violations
from $2$ to $19$ there and from $2$ to $54$ in S1. The policy does
\emph{not} beat the status quo in S1: when the
restart cost is low the winning move is to de-energise aggressively for
long blocks, which the plant already does, and forecast noise costs more
than selectivity gains.

\subsection{Economic sensitivity}
\label{sec:decision:sensitivity}

\paragraph{Break-even is tariff-invariant absent non-energy costs.}
Sweeping the tariff over a $12\times$ range ($0.04$--$0.50$\,USD/kWh)
at $E_r = 10$\,kWh with $\ell = c_{\text{dl}} = 0$ leaves $D^{\star}$
constant at $2.669$\,h, matching $E_r/P_{\text{sb}}$ to seven
significant figures; introducing a $1.00$\,USD non-energy restart cost
makes it span $3.20$--$9.34$\,h. A tariff-versus-break-even map is
therefore uninformative for a purely electrical restart cost, and we
present the plane over restart energy and delay valuation instead, with
tariff as a third swept axis.

\begin{figure*}[t]
\centering
\includegraphics[width=\textwidth]{fig_p7_economic_plane}
\caption{The economic plane at a tariff of $0.12$\,USD/kWh, over the two
coefficients that carry information once break-even is shown to be
tariff-invariant. \textbf{(A)} The derived break-even duration
$D^{\star}$, an output of the optimisation rather than an input
threshold. \textbf{(B)} Annual saving under the proposed policy and
\textbf{(C)} energy saved over the test window, both on a diverging
scale with the zero contour drawn. \textbf{Panels B and C disagree in
sign over part of the plane}: where restarts are expensive the
profitable advice is to stop de-energising, so currency and energy move
in opposite directions and neither can be reported without the other.}
\label{fig:economic-plane}
\end{figure*}

\paragraph{Regimes.}
Figure~\ref{fig:economic-plane} maps the plane.
Over $256$ combinations of tariff, restart energy and delay valuation the
proposed policy is profitable in $249$, but only $171$ also reduce
energy: where restarts are expensive the cheapest advice is to
\emph{stop} de-energising, so currency saved and energy saved carry
opposite signs, and the two must be reported together. The regime
boundaries are
\begin{itemize}
\item the policy overtakes the status quo from $D^{\star} \approx 0.27$\,h;
\item ``never shut down'' becomes optimal from $D^{\star}\approx 27.5$\,h,
      just beyond the longest observed idle episode ($24.1$\,h);
\item the advantage over ski-rental \emph{grows} with $D^{\star}$
      ($+3.1$\,USD in S2, $+17.9$\,USD in S3 over the test window).
\end{itemize}
The last point is the practical message: when waiting is cheap a clock is
almost as good as a forecast, and forecasting earns its place precisely
where restarts are expensive.

\paragraph{Payback.}
At the measured coefficients and $D^{\star}=1$\,h the policy saves
$39$\,USD/yr per machine on the unconstrained labels of
Table~\ref{tab:decision:policies} and $2.9 \pm 0.6$\,USD/yr
(mean $\pm$ sd over ten seeds; Phase~V Task~12C) under the
minimum-dwell labelling, so a $5\,000$\,USD retrofit would pay back in
$130$ years on the first figure and not at all on the second, and a
three-year payback demands a per-machine cost below $116$\,USD in the
most favourable of the two. This follows directly from the plant already de-energising
for the long gaps: only $87$\,h of energised idle remain over $63$ days
of record, some $325$\,kWh. On pelletizer-I the framework's value is
therefore in decision quality --- it captures most of the attainable
head-room while avoiding the $20$ constraint-violating shutdowns the
status quo commits, and, in the headline scenario, its $29$ loss-making
ones --- rather than in a
large absolute saving. Section~\ref{sec:multimachine} examines whether
any other machine in the facility has a larger energised-idle fraction
and so changes this conclusion; none does.

%  Bibliography: paper/references.bib
%  Numbers below are generated by  python run_phase3.py  and stored in
%  outputs/phase3/task6/<machine>/task6_results.json
%             and outputs/phase3/task7/<machine>/task7_results.json
% =====================================================================

\section{Optimisation-Based Shutdown Decision Support}
\label{sec:decision}

\subsection{Problem statement}
\label{sec:decision:problem}

Let $x(t)\in\{0,1\}$ indicate that the machine is energised at time $t$.
Over a horizon the plant incurs standby energy cost while idle and
energised, a restart cost each time the machine is re-energised, and a
production-delay cost because a cold machine reaches production more
slowly than a warm one. The shutdown policy is the solution of

\begin{equation}
\label{eq:objective}
\min_{x(\cdot)}\;
C_{\text{total}}
= c_{\text{sb}}\,h_{\text{standby}}
+ c_{\text{re}}\,N_{\text{restarts}}
+ c_{\text{dl}}\,h_{\text{delay}}
\end{equation}
subject to
\begin{align}
h_{\text{off}}^{(i)} &\ge h_{\min}
&&\text{(motor thermal cool-down)} \label{eq:c1}\\
N_{\text{restarts}}^{(d)} &\le N_{\max}
&&\text{(mechanical wear, per day $d$)} \label{eq:c2}\\
x(t) &\in \{0,1\}. \label{eq:c3}
\end{align}

The coefficients are
$c_{\text{sb}} = \pi\,P_{\text{sb}}$, with $\pi$ the electricity tariff
and $P_{\text{sb}}$ the mean standby power taken from the inferred
STANDBY component of Section~\ref{sec:method}; and
$c_{\text{re}} = \pi E_{r} + \ell + c_{\text{dl}} T_{d}$, with $E_{r}$
the restart energy, $\ell$ any non-energy per-cycle cost, and $T_{d}$
the restart delay. Section~\ref{sec:decision:coeff} measures $P_{\text{sb}}$,
$E_r$ and $T_d$ from the record itself.

\paragraph{The break-even duration is derived, not supplied.}
Setting the saving from de-energising an idle interval of length $D$ to
zero gives
\begin{equation}
\label{eq:breakeven}
D^{\star} \;=\; \frac{c_{\text{re}}}{c_{\text{sb}}}
\;=\; \frac{\pi E_{r} + \ell + c_{\text{dl}} T_{d}}{\pi\,P_{\text{sb}}} .
\end{equation}
This replaces the fixed threshold used in the conventional formulation,
in which $D^{\star}$ is an operator-supplied constant. Equation
\eqref{eq:breakeven} also carries a consequence that is easy to miss:
\emph{when the restart cost is purely electrical} ($\ell = 0$,
$c_{\text{dl}} = 0$) \emph{the tariff cancels}, leaving
$D^{\star} = E_{r}/P_{\text{sb}}$, and the break-even duration is
independent of electricity price. Price enters only through the
non-energy terms. We verify this numerically in
Section~\ref{sec:decision:sensitivity} and it governs how the
sensitivity analysis must be presented.

\subsection{Reduction to one decision per idle episode}
\label{sec:decision:reduction}

Although $x(t)$ is defined at every sample, $C_{\text{total}}$ changes
value only at the boundaries of an \emph{idle episode} --- a maximal
interval of non-productive operation bounded by production. Within one
episode the standby cost is linear and non-decreasing in the time spent
energised, so the saving from de-energising after $\tau$ seconds of an
episode of realised length $D$ is
\begin{equation}
\label{eq:epsave}
S(D,\tau) =
\begin{cases}
c_{\text{sb}}\,(D-\tau) - c_{\text{re}}, & \tau < D,\\
0, & \text{otherwise,}
\end{cases}
\end{equation}
which is decreasing in $\tau$. Any episode that is de-energised at all
is therefore optimally de-energised at its onset, and the problem
reduces to choosing one scalar $\tau_i$ per episode.

\begin{proposition}[Exact offline optimum]
\label{prop:oracle}
With realised durations known, the optimum of
\eqref{eq:objective}--\eqref{eq:c3} is obtained by selecting, within
each day, the $N_{\max}$ episodes of largest positive $S(D_i,0)$ among
those with $D_i \ge h_{\min}$.
\end{proposition}

\noindent
The objective is modular over episodes and the only coupling is the
per-day cardinality constraint \eqref{eq:c2}, so greedy selection is
optimal. Consequently the oracle bound reported below is the exact
solution of the stated problem rather than a heuristic for it, and no
dynamic programme over the $5.5\times10^{6}$ samples is required.

\subsection{Measuring the restart coefficients}
\label{sec:decision:coeff}

Rather than assume $E_r$ and $T_d$, we estimate them from the record.
Every resumption of production is classified by the state the machine
resumed \emph{from} --- cold (OFF) or warm (STANDBY) --- and the cost
attributable to having de-energised is the difference between the two.
Charging the entire cold-start ramp to the shutdown decision would bill
it for a ramp the machine would have paid in any case.

On pelletizer-I the machine reaches 90\% of production power in a median
$574$\,s from cold ($n=126$) against $308$\,s from warm ($n=682$), giving
\begin{equation}
T_{d} = 266\ \text{s}\qquad (\text{95\% CI } [140,\,351]),
\end{equation}
by percentile bootstrap. Energy integrated over a common $1800$\,s
window from the moment production resumes gives
$-6.21$\,kWh (95\% CI $[-7.99,\,-4.64]$; $n=109$ cold, $2526$ warm), and
a milestone-based estimator gives $-1.78$\,kWh.

Both energy estimates are negative and significantly so. This is not an
artefact: a cold start reaches production more slowly, so within any
fixed window it has produced less and therefore drawn less. The energy
difference and the delay are two views of one physical fact, and summing
a negative energy term with a positive delay term would double-count it.
We therefore report that \emph{this machine has no separately
identifiable electrical restart penalty --- its restart cost is a
production-delay cost} --- set $E_r = 0$ in the baseline, and sweep $E_r$
over the range reported in the literature in
Section~\ref{sec:decision:sensitivity}, which also covers machines for
which the penalty is real.

\subsection{Policies}
\label{sec:decision:policies}

Online, $D$ is unknown when the decision must be made. Paying a small
cost repeatedly or a large one-off cost that ends the payments is
exactly the ski-rental (spin-block) structure, and the classical
guarantee applies: de-energising once accumulated standby cost equals
the restart cost is $2$-competitive against the offline optimum for
\emph{any} duration distribution~\cite{karlin1994competitive,karlin1988competitive}.
The same structure underlies online device power-down~\cite{irani2003online}.
This --- not ``never shut down'' --- is the baseline a forecast-driven
policy must beat.

We compare: \textsc{observed}, the plant's own behaviour read off the
inferred labels; \textsc{never} and \textsc{immediate}, the two trivial
policies; \textsc{static}, the conventional rule that forecasts the
duration once at onset and compares it with $D^{\star}$;
\textsc{ski-rental}, the forecast-free elapsed-time rule; the proposed
\textsc{forecast-opt}; and \textsc{oracle} from
Proposition~\ref{prop:oracle}.

\paragraph{Proposed policy.}
Every $60$\,s while the machine remains idle, the rule re-decides using
a forecast of the \emph{remaining} idle time $R$:
\begin{equation}
\label{eq:rule}
\text{de-energise} \iff
\underbrace{c_{\text{sb}}\,\mathbb{E}[R \mid \mathcal{F}_\tau] > c_{\text{re}}}_{\text{economics}}
\;\wedge\;
\underbrace{\Pr(R \ge h_{\min} \mid \mathcal{F}_\tau) \ge 1-\alpha}_{\text{feasibility}} ,
\end{equation}
where $\mathcal{F}_\tau$ is the information available at elapsed time
$\tau$: elapsed idle time, the \emph{current} position in the factory
calendar, and the load history of the production run that just ended.

Two aspects of \eqref{eq:rule} are deliberate. First, it is
\emph{sequential} rather than a one-shot classification at onset --- the
moment when the least information is available --- so each additional
idle minute is used as evidence. Second, the two conditions are not
redundant. Because $S$ is affine in $R$ by \eqref{eq:epsave}, the
unconstrained risk-neutral rule depends on the predictive distribution
only through its \emph{mean}, and a mean regression suffices. Constraint
\eqref{eq:c1} breaks that affinity: a shutdown whose idle period proves
shorter than $h_{\min}$ is infeasible rather than merely unprofitable,
and no expectation over $R$ expresses this. Feasibility is therefore
imposed as a chance constraint~\cite{charnes1959chance} at a stated
confidence $1-\alpha = 0.90$, using a separate classifier for
$\Pr(R\ge h_{\min})$ (AUC $0.939$, Brier $0.102$). Setting $\alpha=1$
recovers the pure risk-neutral rule and isolates the constraint's
contribution.

\paragraph{The estimand matters more than the model.}
Idle durations span four orders of magnitude, so fitting
$\log(1+R)$ is the natural choice --- and it is the wrong one here.
Back-transforming a log-scale fit estimates the conditional
\emph{median}, and Duan's smearing correction~\cite{duan1983smearing}
repairs it with a single global constant that cannot absorb a
covariate-dependent bias. At episode onset the conditional median is
$7$\,s while the conditional mean exceeds $2\,000$\,s. Since
\eqref{eq:rule} consumes the mean, the estimand must be elicited by the
loss used to fit it~\cite{gneiting2011making}; we therefore fit the raw
scale under squared error and retain the log-scale model as an ablation.

\subsection{Experimental protocol}
\label{sec:decision:protocol}

Spike rows are retained: removing $19.5\%$ of a $1$\,Hz series destroys
the correspondence between row counts and seconds on which every episode
duration depends. Episodes touching one of the record's $31$
acquisition gaps are excluded ($51$ episodes, $187.9$\,h), because their
true extent is unknown and counting them would manufacture savings out
of missing data. This leaves $3\,224$ usable episodes spanning
$474.9$\,h of idle time. The split is \emph{chronological} ($70/30$):
unlike the unsupervised comparison of Section~\ref{sec:results}, a
decision policy is deployed forward in time, and a shuffled split would
let it learn the future duty cycle. The per-day restart cap is enforced
at evaluation time, identically for every policy. All annualisation uses
the $63.4$ days the record actually covers, not its $154$-day span.

\paragraph{The status quo is the baseline.}
Of the $474.9$ usable idle hours the plant already de-energises for
$388.1$, leaving $86.8$\,h energised. Savings quoted against a
``never shut down'' hypothetical would therefore be savings the plant has
already banked; every figure below is quoted against \textsc{observed}.

\paragraph{Scenarios.}
Every coefficient is measurable except the valuation of lost production,
and $D^{\star}$ is linear in it, so asserting a value would silently
determine the conclusion. Scenarios are instead indexed by where
$D^{\star}$ falls within the observed idle-duration distribution, and the
delay valuation each implies is reported as an \emph{output}
(Table~\ref{tab:decision:scenarios}).

\begin{table}[t]
\centering
\caption{Economic scenarios, indexed by the derived break-even point.
The implied production-delay valuation is an output of
Eq.~\eqref{eq:breakeven}, not an assumption.}
\label{tab:decision:scenarios}
\begin{tabular}{lrrrr}
\hline
Scenario & $D^{\star}$ & Implied $c_{\text{dl}}$ & $c_{\text{re}}$ & Episodes above \\
         & (min) & (USD/h) & (USD) & \\
\hline
S1 low       &  10 &  1.01 & 0.075 & 228 \\
S2 moderate  &  60 &  6.09 & 0.450 &  53 \\
S3 high      & 240 & 24.34 & 1.799 &  43 \\
\hline
\end{tabular}
\end{table}

\subsection{Results}
\label{sec:decision:results}

Table~\ref{tab:decision:policies} reports annual saving against the
status quo and the share of the oracle head-room captured, on $968$
held-out episodes spanning $78.3$ days.

\begin{table}[t]
\centering
\caption{Policy comparison on held-out episodes: annual saving versus
the plant's own behaviour (USD/yr) and share of the attainable head-room
(\%). The three middle rows are ablations of the proposed policy.}
\label{tab:decision:policies}
\begin{tabular}{lrrrrrr}
\hline
& \multicolumn{2}{c}{S1 (10\,min)} & \multicolumn{2}{c}{S2 (60\,min)} & \multicolumn{2}{c}{S3 (240\,min)}\\
Policy & USD/yr & \% & USD/yr & \% & USD/yr & \% \\
\hline
Never shut down            & $-215$ & $-530$ & $-133$ & $-190$ & $+163$ & $65$ \\
Observed (status quo)      & $0$ & $0$ & $0$ & $0$ & $0$ & $0$ \\
Immediate shutdown         & $-237$ & $-584$ & $-316$ & $-451$ & $-599$ & $-241$ \\
Static break-even          & $-98$ & $-242$ & $-46$ & $-65$ & $+175$ & $70$ \\
Ski-rental (no forecast)   & $-23$ & $-56$ & $+24$ & $34$ & $+106$ & $43$ \\
\quad ablation: median estimand      & $-23$ & $-56$ & $-14$ & $-20$ & $+180$ & $72$ \\
\quad ablation: no chance constraint & $-73$ & $-180$ & $+6$ & $8$ & $+173$ & $69$ \\
\textbf{Forecast optimisation}       & $\mathbf{-3}$ & $\mathbf{-8}$ & $\mathbf{+39}$ & $\mathbf{55}$ & $\mathbf{+190}$ & $\mathbf{76}$ \\
Oracle (offline optimum)   & $+41$ & $100$ & $+70$ & $100$ & $+249$ & $100$ \\
\hline
\end{tabular}
\end{table}

\paragraph{These figures are conditional on the labelling.}
Table~\ref{tab:decision:policies} is computed on the state labels of
Section~\ref{sec:states}, in which the Viterbi path is left
unconstrained. Re-deriving the same table from labels carrying the
minimum-dwell constraint of Section~\ref{sec:flicker}---the same
machine, the same policy, the same coefficients, and a test window cut
at the same instant---changes the S2 column materially: the attainable
head-room falls from $71$\,USD/yr to $38$\,USD/yr and the proposed
policy from $+41$\,USD/yr ($57\%$ of the oracle) to
$+2.9 \pm 0.6$\,USD/yr ($7.8 \pm 1.6\%$; mean and standard deviation
over ten random seeds from Phase~V Task~12C). The physical quantities
are unmoved (idle $474.9$ vs $474.1$\,h; energised idle $86.8$ vs
$87.8$\,h; standby power $3\,747$ vs $4\,208$\,W); what moves is the
\emph{partition of idle time into episodes}, which is the unit the
decision layer bids on: enforcing a minimum dwell merges $3\,224$
episodes into $1\,312$ and raises the median episode from $12$ to
$90$\,s. Sub-ten-second ``opportunities'' created by label flicker are
not shutdown opportunities, so the smaller figure is the defensible
one, and both are reported here rather than the more favourable one
alone. The seed-to-seed standard deviation ($0.6$\,USD/yr) is small
relative to the bootstrap sampling interval ($\pm 50$\,USD/yr over
18 factory days; Phase~V Task~11B), confirming that model-fitting
variance does not explain the uncertainty.

The proposed policy captures $55$--$76\%$ of the attainable head-room in
S2 and S3 and beats the forecast-free ski-rental rule in all three
scenarios. Both ablations are material: replacing the mean estimand with
the log-scale median costs $53$\,USD/yr in S2, and removing the chance
constraint costs $33$\,USD/yr in S2 while raising minimum-off violations
from $2$ to $19$ there and from $2$ to $54$ in S1. The policy does
\emph{not} beat the status quo in S1: when the
restart cost is low the winning move is to de-energise aggressively for
long blocks, which the plant already does, and forecast noise costs more
than selectivity gains.

\subsection{Economic sensitivity}
\label{sec:decision:sensitivity}

\paragraph{Break-even is tariff-invariant absent non-energy costs.}
Sweeping the tariff over a $12\times$ range ($0.04$--$0.50$\,USD/kWh)
at $E_r = 10$\,kWh with $\ell = c_{\text{dl}} = 0$ leaves $D^{\star}$
constant at $2.669$\,h, matching $E_r/P_{\text{sb}}$ to seven
significant figures; introducing a $1.00$\,USD non-energy restart cost
makes it span $3.20$--$9.34$\,h. A tariff-versus-break-even map is
therefore uninformative for a purely electrical restart cost, and we
present the plane over restart energy and delay valuation instead, with
tariff as a third swept axis.

\begin{figure*}[t]
\centering
\includegraphics[width=\textwidth]{fig_p7_economic_plane}
\caption{The economic plane at a tariff of $0.12$\,USD/kWh, over the two
coefficients that carry information once break-even is shown to be
tariff-invariant. \textbf{(A)} The derived break-even duration
$D^{\star}$, an output of the optimisation rather than an input
threshold. \textbf{(B)} Annual saving under the proposed policy and
\textbf{(C)} energy saved over the test window, both on a diverging
scale with the zero contour drawn. \textbf{Panels B and C disagree in
sign over part of the plane}: where restarts are expensive the
profitable advice is to stop de-energising, so currency and energy move
in opposite directions and neither can be reported without the other.}
\label{fig:economic-plane}
\end{figure*}

\paragraph{Regimes.}
Figure~\ref{fig:economic-plane} maps the plane.
Over $256$ combinations of tariff, restart energy and delay valuation the
proposed policy is profitable in $249$, but only $171$ also reduce
energy: where restarts are expensive the cheapest advice is to
\emph{stop} de-energising, so currency saved and energy saved carry
opposite signs, and the two must be reported together. The regime
boundaries are
\begin{itemize}
\item the policy overtakes the status quo from $D^{\star} \approx 0.27$\,h;
\item ``never shut down'' becomes optimal from $D^{\star}\approx 27.5$\,h,
      just beyond the longest observed idle episode ($24.1$\,h);
\item the advantage over ski-rental \emph{grows} with $D^{\star}$
      ($+3.1$\,USD in S2, $+17.9$\,USD in S3 over the test window).
\end{itemize}
The last point is the practical message: when waiting is cheap a clock is
almost as good as a forecast, and forecasting earns its place precisely
where restarts are expensive.

\paragraph{Payback.}
At the measured coefficients and $D^{\star}=1$\,h the policy saves
$39$\,USD/yr per machine on the unconstrained labels of
Table~\ref{tab:decision:policies} and $2.9 \pm 0.6$\,USD/yr
(mean $\pm$ sd over ten seeds; Phase~V Task~12C) under the
minimum-dwell labelling, so a $5\,000$\,USD retrofit would pay back in
$130$ years on the first figure and not at all on the second, and a
three-year payback demands a per-machine cost below $116$\,USD in the
most favourable of the two. This follows directly from the plant already de-energising
for the long gaps: only $87$\,h of energised idle remain over $63$ days
of record, some $325$\,kWh. On pelletizer-I the framework's value is
therefore in decision quality --- it captures most of the attainable
head-room while avoiding the $20$ constraint-violating shutdowns the
status quo commits, and, in the headline scenario, its $29$ loss-making
ones --- rather than in a
large absolute saving. Section~\ref{sec:multimachine} examines whether
any other machine in the facility has a larger energised-idle fraction
and so changes this conclusion; none does.

%  Bibliography: paper/references.bib
%  Numbers below are generated by  python run_phase3.py  and stored in
%  outputs/phase3/task6/<machine>/task6_results.json
%             and outputs/phase3/task7/<machine>/task7_results.json
% =====================================================================

\section{Optimisation-Based Shutdown Decision Support}
\label{sec:decision}

\subsection{Problem statement}
\label{sec:decision:problem}

Let $x(t)\in\{0,1\}$ indicate that the machine is energised at time $t$.
Over a horizon the plant incurs standby energy cost while idle and
energised, a restart cost each time the machine is re-energised, and a
production-delay cost because a cold machine reaches production more
slowly than a warm one. The shutdown policy is the solution of

\begin{equation}
\label{eq:objective}
\min_{x(\cdot)}\;
C_{\text{total}}
= c_{\text{sb}}\,h_{\text{standby}}
+ c_{\text{re}}\,N_{\text{restarts}}
+ c_{\text{dl}}\,h_{\text{delay}}
\end{equation}
subject to
\begin{align}
h_{\text{off}}^{(i)} &\ge h_{\min}
&&\text{(motor thermal cool-down)} \label{eq:c1}\\
N_{\text{restarts}}^{(d)} &\le N_{\max}
&&\text{(mechanical wear, per day $d$)} \label{eq:c2}\\
x(t) &\in \{0,1\}. \label{eq:c3}
\end{align}

The coefficients are
$c_{\text{sb}} = \pi\,P_{\text{sb}}$, with $\pi$ the electricity tariff
and $P_{\text{sb}}$ the mean standby power taken from the inferred
STANDBY component of Section~\ref{sec:method}; and
$c_{\text{re}} = \pi E_{r} + \ell + c_{\text{dl}} T_{d}$, with $E_{r}$
the restart energy, $\ell$ any non-energy per-cycle cost, and $T_{d}$
the restart delay. Section~\ref{sec:decision:coeff} measures $P_{\text{sb}}$,
$E_r$ and $T_d$ from the record itself.

\paragraph{The break-even duration is derived, not supplied.}
Setting the saving from de-energising an idle interval of length $D$ to
zero gives
\begin{equation}
\label{eq:breakeven}
D^{\star} \;=\; \frac{c_{\text{re}}}{c_{\text{sb}}}
\;=\; \frac{\pi E_{r} + \ell + c_{\text{dl}} T_{d}}{\pi\,P_{\text{sb}}} .
\end{equation}
This replaces the fixed threshold used in the conventional formulation,
in which $D^{\star}$ is an operator-supplied constant. Equation
\eqref{eq:breakeven} also carries a consequence that is easy to miss:
\emph{when the restart cost is purely electrical} ($\ell = 0$,
$c_{\text{dl}} = 0$) \emph{the tariff cancels}, leaving
$D^{\star} = E_{r}/P_{\text{sb}}$, and the break-even duration is
independent of electricity price. Price enters only through the
non-energy terms. We verify this numerically in
Section~\ref{sec:decision:sensitivity} and it governs how the
sensitivity analysis must be presented.

\subsection{Reduction to one decision per idle episode}
\label{sec:decision:reduction}

Although $x(t)$ is defined at every sample, $C_{\text{total}}$ changes
value only at the boundaries of an \emph{idle episode} --- a maximal
interval of non-productive operation bounded by production. Within one
episode the standby cost is linear and non-decreasing in the time spent
energised, so the saving from de-energising after $\tau$ seconds of an
episode of realised length $D$ is
\begin{equation}
\label{eq:epsave}
S(D,\tau) =
\begin{cases}
c_{\text{sb}}\,(D-\tau) - c_{\text{re}}, & \tau < D,\\
0, & \text{otherwise,}
\end{cases}
\end{equation}
which is decreasing in $\tau$. Any episode that is de-energised at all
is therefore optimally de-energised at its onset, and the problem
reduces to choosing one scalar $\tau_i$ per episode.

\begin{proposition}[Exact offline optimum]
\label{prop:oracle}
With realised durations known, the optimum of
\eqref{eq:objective}--\eqref{eq:c3} is obtained by selecting, within
each day, the $N_{\max}$ episodes of largest positive $S(D_i,0)$ among
those with $D_i \ge h_{\min}$.
\end{proposition}

\noindent
The objective is modular over episodes and the only coupling is the
per-day cardinality constraint \eqref{eq:c2}, so greedy selection is
optimal. Consequently the oracle bound reported below is the exact
solution of the stated problem rather than a heuristic for it, and no
dynamic programme over the $5.5\times10^{6}$ samples is required.

\subsection{Measuring the restart coefficients}
\label{sec:decision:coeff}

Rather than assume $E_r$ and $T_d$, we estimate them from the record.
Every resumption of production is classified by the state the machine
resumed \emph{from} --- cold (OFF) or warm (STANDBY) --- and the cost
attributable to having de-energised is the difference between the two.
Charging the entire cold-start ramp to the shutdown decision would bill
it for a ramp the machine would have paid in any case.

On pelletizer-I the machine reaches 90\% of production power in a median
$574$\,s from cold ($n=126$) against $308$\,s from warm ($n=682$), giving
\begin{equation}
T_{d} = 266\ \text{s}\qquad (\text{95\% CI } [140,\,351]),
\end{equation}
by percentile bootstrap. Energy integrated over a common $1800$\,s
window from the moment production resumes gives
$-6.21$\,kWh (95\% CI $[-7.99,\,-4.64]$; $n=109$ cold, $2526$ warm), and
a milestone-based estimator gives $-1.78$\,kWh.

Both energy estimates are negative and significantly so. This is not an
artefact: a cold start reaches production more slowly, so within any
fixed window it has produced less and therefore drawn less. The energy
difference and the delay are two views of one physical fact, and summing
a negative energy term with a positive delay term would double-count it.
We therefore report that \emph{this machine has no separately
identifiable electrical restart penalty --- its restart cost is a
production-delay cost} --- set $E_r = 0$ in the baseline, and sweep $E_r$
over the range reported in the literature in
Section~\ref{sec:decision:sensitivity}, which also covers machines for
which the penalty is real.

\subsection{Policies}
\label{sec:decision:policies}

Online, $D$ is unknown when the decision must be made. Paying a small
cost repeatedly or a large one-off cost that ends the payments is
exactly the ski-rental (spin-block) structure, and the classical
guarantee applies: de-energising once accumulated standby cost equals
the restart cost is $2$-competitive against the offline optimum for
\emph{any} duration distribution~\cite{karlin1994competitive,karlin1988competitive}.
The same structure underlies online device power-down~\cite{irani2003online}.
This --- not ``never shut down'' --- is the baseline a forecast-driven
policy must beat.

We compare: \textsc{observed}, the plant's own behaviour read off the
inferred labels; \textsc{never} and \textsc{immediate}, the two trivial
policies; \textsc{static}, the conventional rule that forecasts the
duration once at onset and compares it with $D^{\star}$;
\textsc{ski-rental}, the forecast-free elapsed-time rule; the proposed
\textsc{forecast-opt}; and \textsc{oracle} from
Proposition~\ref{prop:oracle}.

\paragraph{Proposed policy.}
Every $60$\,s while the machine remains idle, the rule re-decides using
a forecast of the \emph{remaining} idle time $R$:
\begin{equation}
\label{eq:rule}
\text{de-energise} \iff
\underbrace{c_{\text{sb}}\,\mathbb{E}[R \mid \mathcal{F}_\tau] > c_{\text{re}}}_{\text{economics}}
\;\wedge\;
\underbrace{\Pr(R \ge h_{\min} \mid \mathcal{F}_\tau) \ge 1-\alpha}_{\text{feasibility}} ,
\end{equation}
where $\mathcal{F}_\tau$ is the information available at elapsed time
$\tau$: elapsed idle time, the \emph{current} position in the factory
calendar, and the load history of the production run that just ended.

Two aspects of \eqref{eq:rule} are deliberate. First, it is
\emph{sequential} rather than a one-shot classification at onset --- the
moment when the least information is available --- so each additional
idle minute is used as evidence. Second, the two conditions are not
redundant. Because $S$ is affine in $R$ by \eqref{eq:epsave}, the
unconstrained risk-neutral rule depends on the predictive distribution
only through its \emph{mean}, and a mean regression suffices. Constraint
\eqref{eq:c1} breaks that affinity: a shutdown whose idle period proves
shorter than $h_{\min}$ is infeasible rather than merely unprofitable,
and no expectation over $R$ expresses this. Feasibility is therefore
imposed as a chance constraint~\cite{charnes1959chance} at a stated
confidence $1-\alpha = 0.90$, using a separate classifier for
$\Pr(R\ge h_{\min})$ (AUC $0.939$, Brier $0.102$). Setting $\alpha=1$
recovers the pure risk-neutral rule and isolates the constraint's
contribution.

\paragraph{The estimand matters more than the model.}
Idle durations span four orders of magnitude, so fitting
$\log(1+R)$ is the natural choice --- and it is the wrong one here.
Back-transforming a log-scale fit estimates the conditional
\emph{median}, and Duan's smearing correction~\cite{duan1983smearing}
repairs it with a single global constant that cannot absorb a
covariate-dependent bias. At episode onset the conditional median is
$7$\,s while the conditional mean exceeds $2\,000$\,s. Since
\eqref{eq:rule} consumes the mean, the estimand must be elicited by the
loss used to fit it~\cite{gneiting2011making}; we therefore fit the raw
scale under squared error and retain the log-scale model as an ablation.

\subsection{Experimental protocol}
\label{sec:decision:protocol}

Spike rows are retained: removing $19.5\%$ of a $1$\,Hz series destroys
the correspondence between row counts and seconds on which every episode
duration depends. Episodes touching one of the record's $31$
acquisition gaps are excluded ($51$ episodes, $187.9$\,h), because their
true extent is unknown and counting them would manufacture savings out
of missing data. This leaves $3\,224$ usable episodes spanning
$474.9$\,h of idle time. The split is \emph{chronological} ($70/30$):
unlike the unsupervised comparison of Section~\ref{sec:results}, a
decision policy is deployed forward in time, and a shuffled split would
let it learn the future duty cycle. The per-day restart cap is enforced
at evaluation time, identically for every policy. All annualisation uses
the $63.4$ days the record actually covers, not its $154$-day span.

\paragraph{The status quo is the baseline.}
Of the $474.9$ usable idle hours the plant already de-energises for
$388.1$, leaving $86.8$\,h energised. Savings quoted against a
``never shut down'' hypothetical would therefore be savings the plant has
already banked; every figure below is quoted against \textsc{observed}.

\paragraph{Scenarios.}
Every coefficient is measurable except the valuation of lost production,
and $D^{\star}$ is linear in it, so asserting a value would silently
determine the conclusion. Scenarios are instead indexed by where
$D^{\star}$ falls within the observed idle-duration distribution, and the
delay valuation each implies is reported as an \emph{output}
(Table~\ref{tab:decision:scenarios}).

\begin{table}[t]
\centering
\caption{Economic scenarios, indexed by the derived break-even point.
The implied production-delay valuation is an output of
Eq.~\eqref{eq:breakeven}, not an assumption.}
\label{tab:decision:scenarios}
\begin{tabular}{lrrrr}
\hline
Scenario & $D^{\star}$ & Implied $c_{\text{dl}}$ & $c_{\text{re}}$ & Episodes above \\
         & (min) & (USD/h) & (USD) & \\
\hline
S1 low       &  10 &  1.01 & 0.075 & 228 \\
S2 moderate  &  60 &  6.09 & 0.450 &  53 \\
S3 high      & 240 & 24.34 & 1.799 &  43 \\
\hline
\end{tabular}
\end{table}

\subsection{Results}
\label{sec:decision:results}

Table~\ref{tab:decision:policies} reports annual saving against the
status quo and the share of the oracle head-room captured, on $968$
held-out episodes spanning $78.3$ days.

\begin{table}[t]
\centering
\caption{Policy comparison on held-out episodes: annual saving versus
the plant's own behaviour (USD/yr) and share of the attainable head-room
(\%). The three middle rows are ablations of the proposed policy.}
\label{tab:decision:policies}
\begin{tabular}{lrrrrrr}
\hline
& \multicolumn{2}{c}{S1 (10\,min)} & \multicolumn{2}{c}{S2 (60\,min)} & \multicolumn{2}{c}{S3 (240\,min)}\\
Policy & USD/yr & \% & USD/yr & \% & USD/yr & \% \\
\hline
Never shut down            & $-215$ & $-530$ & $-133$ & $-190$ & $+163$ & $65$ \\
Observed (status quo)      & $0$ & $0$ & $0$ & $0$ & $0$ & $0$ \\
Immediate shutdown         & $-237$ & $-584$ & $-316$ & $-451$ & $-599$ & $-241$ \\
Static break-even          & $-98$ & $-242$ & $-46$ & $-65$ & $+175$ & $70$ \\
Ski-rental (no forecast)   & $-23$ & $-56$ & $+24$ & $34$ & $+106$ & $43$ \\
\quad ablation: median estimand      & $-23$ & $-56$ & $-14$ & $-20$ & $+180$ & $72$ \\
\quad ablation: no chance constraint & $-73$ & $-180$ & $+6$ & $8$ & $+173$ & $69$ \\
\textbf{Forecast optimisation}       & $\mathbf{-3}$ & $\mathbf{-8}$ & $\mathbf{+39}$ & $\mathbf{55}$ & $\mathbf{+190}$ & $\mathbf{76}$ \\
Oracle (offline optimum)   & $+41$ & $100$ & $+70$ & $100$ & $+249$ & $100$ \\
\hline
\end{tabular}
\end{table}

\paragraph{These figures are conditional on the labelling.}
Table~\ref{tab:decision:policies} is computed on the state labels of
Section~\ref{sec:states}, in which the Viterbi path is left
unconstrained. Re-deriving the same table from labels carrying the
minimum-dwell constraint of Section~\ref{sec:flicker}---the same
machine, the same policy, the same coefficients, and a test window cut
at the same instant---changes the S2 column materially: the attainable
head-room falls from $71$\,USD/yr to $38$\,USD/yr and the proposed
policy from $+41$\,USD/yr ($57\%$ of the oracle) to
$+2.9 \pm 0.6$\,USD/yr ($7.8 \pm 1.6\%$; mean and standard deviation
over ten random seeds from Phase~V Task~12C). The physical quantities
are unmoved (idle $474.9$ vs $474.1$\,h; energised idle $86.8$ vs
$87.8$\,h; standby power $3\,747$ vs $4\,208$\,W); what moves is the
\emph{partition of idle time into episodes}, which is the unit the
decision layer bids on: enforcing a minimum dwell merges $3\,224$
episodes into $1\,312$ and raises the median episode from $12$ to
$90$\,s. Sub-ten-second ``opportunities'' created by label flicker are
not shutdown opportunities, so the smaller figure is the defensible
one, and both are reported here rather than the more favourable one
alone. The seed-to-seed standard deviation ($0.6$\,USD/yr) is small
relative to the bootstrap sampling interval ($\pm 50$\,USD/yr over
18 factory days; Phase~V Task~11B), confirming that model-fitting
variance does not explain the uncertainty.

The proposed policy captures $55$--$76\%$ of the attainable head-room in
S2 and S3 and beats the forecast-free ski-rental rule in all three
scenarios. Both ablations are material: replacing the mean estimand with
the log-scale median costs $53$\,USD/yr in S2, and removing the chance
constraint costs $33$\,USD/yr in S2 while raising minimum-off violations
from $2$ to $19$ there and from $2$ to $54$ in S1. The policy does
\emph{not} beat the status quo in S1: when the
restart cost is low the winning move is to de-energise aggressively for
long blocks, which the plant already does, and forecast noise costs more
than selectivity gains.

\subsection{Economic sensitivity}
\label{sec:decision:sensitivity}

\paragraph{Break-even is tariff-invariant absent non-energy costs.}
Sweeping the tariff over a $12\times$ range ($0.04$--$0.50$\,USD/kWh)
at $E_r = 10$\,kWh with $\ell = c_{\text{dl}} = 0$ leaves $D^{\star}$
constant at $2.669$\,h, matching $E_r/P_{\text{sb}}$ to seven
significant figures; introducing a $1.00$\,USD non-energy restart cost
makes it span $3.20$--$9.34$\,h. A tariff-versus-break-even map is
therefore uninformative for a purely electrical restart cost, and we
present the plane over restart energy and delay valuation instead, with
tariff as a third swept axis.

\begin{figure*}[t]
\centering
\includegraphics[width=\textwidth]{fig_p7_economic_plane}
\caption{The economic plane at a tariff of $0.12$\,USD/kWh, over the two
coefficients that carry information once break-even is shown to be
tariff-invariant. \textbf{(A)} The derived break-even duration
$D^{\star}$, an output of the optimisation rather than an input
threshold. \textbf{(B)} Annual saving under the proposed policy and
\textbf{(C)} energy saved over the test window, both on a diverging
scale with the zero contour drawn. \textbf{Panels B and C disagree in
sign over part of the plane}: where restarts are expensive the
profitable advice is to stop de-energising, so currency and energy move
in opposite directions and neither can be reported without the other.}
\label{fig:economic-plane}
\end{figure*}

\paragraph{Regimes.}
Figure~\ref{fig:economic-plane} maps the plane.
Over $256$ combinations of tariff, restart energy and delay valuation the
proposed policy is profitable in $249$, but only $171$ also reduce
energy: where restarts are expensive the cheapest advice is to
\emph{stop} de-energising, so currency saved and energy saved carry
opposite signs, and the two must be reported together. The regime
boundaries are
\begin{itemize}
\item the policy overtakes the status quo from $D^{\star} \approx 0.27$\,h;
\item ``never shut down'' becomes optimal from $D^{\star}\approx 27.5$\,h,
      just beyond the longest observed idle episode ($24.1$\,h);
\item the advantage over ski-rental \emph{grows} with $D^{\star}$
      ($+3.1$\,USD in S2, $+17.9$\,USD in S3 over the test window).
\end{itemize}
The last point is the practical message: when waiting is cheap a clock is
almost as good as a forecast, and forecasting earns its place precisely
where restarts are expensive.

\paragraph{Payback.}
At the measured coefficients and $D^{\star}=1$\,h the policy saves
$39$\,USD/yr per machine on the unconstrained labels of
Table~\ref{tab:decision:policies} and $2.9 \pm 0.6$\,USD/yr
(mean $\pm$ sd over ten seeds; Phase~V Task~12C) under the
minimum-dwell labelling, so a $5\,000$\,USD retrofit would pay back in
$130$ years on the first figure and not at all on the second, and a
three-year payback demands a per-machine cost below $116$\,USD in the
most favourable of the two. This follows directly from the plant already de-energising
for the long gaps: only $87$\,h of energised idle remain over $63$ days
of record, some $325$\,kWh. On pelletizer-I the framework's value is
therefore in decision quality --- it captures most of the attainable
head-room while avoiding the $20$ constraint-violating shutdowns the
status quo commits, and, in the headline scenario, its $29$ loss-making
ones --- rather than in a
large absolute saving. Section~\ref{sec:multimachine} examines whether
any other machine in the facility has a larger energised-idle fraction
and so changes this conclusion; none does.

%  Bibliography: paper/references.bib
%  Numbers below are generated by  python run_phase3.py  and stored in
%  outputs/phase3/task6/<machine>/task6_results.json
%             and outputs/phase3/task7/<machine>/task7_results.json
% =====================================================================

\section{Optimisation-Based Shutdown Decision Support}
\label{sec:decision}

\subsection{Problem statement}
\label{sec:decision:problem}

Let $x(t)\in\{0,1\}$ indicate that the machine is energised at time $t$.
Over a horizon the plant incurs standby energy cost while idle and
energised, a restart cost each time the machine is re-energised, and a
production-delay cost because a cold machine reaches production more
slowly than a warm one. The shutdown policy is the solution of

\begin{equation}
\label{eq:objective}
\min_{x(\cdot)}\;
C_{\text{total}}
= c_{\text{sb}}\,h_{\text{standby}}
+ c_{\text{re}}\,N_{\text{restarts}}
+ c_{\text{dl}}\,h_{\text{delay}}
\end{equation}
subject to
\begin{align}
h_{\text{off}}^{(i)} &\ge h_{\min}
&&\text{(motor thermal cool-down)} \label{eq:c1}\\
N_{\text{restarts}}^{(d)} &\le N_{\max}
&&\text{(mechanical wear, per day $d$)} \label{eq:c2}\\
x(t) &\in \{0,1\}. \label{eq:c3}
\end{align}

The coefficients are
$c_{\text{sb}} = \pi\,P_{\text{sb}}$, with $\pi$ the electricity tariff
and $P_{\text{sb}}$ the mean standby power taken from the inferred
STANDBY component of Section~\ref{sec:method}; and
$c_{\text{re}} = \pi E_{r} + \ell + c_{\text{dl}} T_{d}$, with $E_{r}$
the restart energy, $\ell$ any non-energy per-cycle cost, and $T_{d}$
the restart delay. Section~\ref{sec:decision:coeff} measures $P_{\text{sb}}$,
$E_r$ and $T_d$ from the record itself.

\paragraph{The break-even duration is derived, not supplied.}
Setting the saving from de-energising an idle interval of length $D$ to
zero gives
\begin{equation}
\label{eq:breakeven}
D^{\star} \;=\; \frac{c_{\text{re}}}{c_{\text{sb}}}
\;=\; \frac{\pi E_{r} + \ell + c_{\text{dl}} T_{d}}{\pi\,P_{\text{sb}}} .
\end{equation}
This replaces the fixed threshold used in the conventional formulation,
in which $D^{\star}$ is an operator-supplied constant. Equation
\eqref{eq:breakeven} also carries a consequence that is easy to miss:
\emph{when the restart cost is purely electrical} ($\ell = 0$,
$c_{\text{dl}} = 0$) \emph{the tariff cancels}, leaving
$D^{\star} = E_{r}/P_{\text{sb}}$, and the break-even duration is
independent of electricity price. Price enters only through the
non-energy terms. We verify this numerically in
Section~\ref{sec:decision:sensitivity} and it governs how the
sensitivity analysis must be presented.

\subsection{Reduction to one decision per idle episode}
\label{sec:decision:reduction}

Although $x(t)$ is defined at every sample, $C_{\text{total}}$ changes
value only at the boundaries of an \emph{idle episode} --- a maximal
interval of non-productive operation bounded by production. Within one
episode the standby cost is linear and non-decreasing in the time spent
energised, so the saving from de-energising after $\tau$ seconds of an
episode of realised length $D$ is
\begin{equation}
\label{eq:epsave}
S(D,\tau) =
\begin{cases}
c_{\text{sb}}\,(D-\tau) - c_{\text{re}}, & \tau < D,\\
0, & \text{otherwise,}
\end{cases}
\end{equation}
which is decreasing in $\tau$. Any episode that is de-energised at all
is therefore optimally de-energised at its onset, and the problem
reduces to choosing one scalar $\tau_i$ per episode.

\begin{proposition}[Exact offline optimum]
\label{prop:oracle}
With realised durations known, the optimum of
\eqref{eq:objective}--\eqref{eq:c3} is obtained by selecting, within
each day, the $N_{\max}$ episodes of largest positive $S(D_i,0)$ among
those with $D_i \ge h_{\min}$.
\end{proposition}

\noindent
The objective is modular over episodes and the only coupling is the
per-day cardinality constraint \eqref{eq:c2}, so greedy selection is
optimal. Consequently the oracle bound reported below is the exact
solution of the stated problem rather than a heuristic for it, and no
dynamic programme over the $5.5\times10^{6}$ samples is required.

\subsection{Measuring the restart coefficients}
\label{sec:decision:coeff}

Rather than assume $E_r$ and $T_d$, we estimate them from the record.
Every resumption of production is classified by the state the machine
resumed \emph{from} --- cold (OFF) or warm (STANDBY) --- and the cost
attributable to having de-energised is the difference between the two.
Charging the entire cold-start ramp to the shutdown decision would bill
it for a ramp the machine would have paid in any case.

On pelletizer-I the machine reaches 90\% of production power in a median
$574$\,s from cold ($n=126$) against $308$\,s from warm ($n=682$), giving
\begin{equation}
T_{d} = 266\ \text{s}\qquad (\text{95\% CI } [140,\,351]),
\end{equation}
by percentile bootstrap. Energy integrated over a common $1800$\,s
window from the moment production resumes gives
$-6.21$\,kWh (95\% CI $[-7.99,\,-4.64]$; $n=109$ cold, $2526$ warm), and
a milestone-based estimator gives $-1.78$\,kWh.

Both energy estimates are negative and significantly so. This is not an
artefact: a cold start reaches production more slowly, so within any
fixed window it has produced less and therefore drawn less. The energy
difference and the delay are two views of one physical fact, and summing
a negative energy term with a positive delay term would double-count it.
We therefore report that \emph{this machine has no separately
identifiable electrical restart penalty --- its restart cost is a
production-delay cost} --- set $E_r = 0$ in the baseline, and sweep $E_r$
over the range reported in the literature in
Section~\ref{sec:decision:sensitivity}, which also covers machines for
which the penalty is real.

\subsection{Policies}
\label{sec:decision:policies}

Online, $D$ is unknown when the decision must be made. Paying a small
cost repeatedly or a large one-off cost that ends the payments is
exactly the ski-rental (spin-block) structure, and the classical
guarantee applies: de-energising once accumulated standby cost equals
the restart cost is $2$-competitive against the offline optimum for
\emph{any} duration distribution~\cite{karlin1994competitive,karlin1988competitive}.
The same structure underlies online device power-down~\cite{irani2003online}.
This --- not ``never shut down'' --- is the baseline a forecast-driven
policy must beat.

We compare: \textsc{observed}, the plant's own behaviour read off the
inferred labels; \textsc{never} and \textsc{immediate}, the two trivial
policies; \textsc{static}, the conventional rule that forecasts the
duration once at onset and compares it with $D^{\star}$;
\textsc{ski-rental}, the forecast-free elapsed-time rule; the proposed
\textsc{forecast-opt}; and \textsc{oracle} from
Proposition~\ref{prop:oracle}.

\paragraph{Proposed policy.}
Every $60$\,s while the machine remains idle, the rule re-decides using
a forecast of the \emph{remaining} idle time $R$:
\begin{equation}
\label{eq:rule}
\text{de-energise} \iff
\underbrace{c_{\text{sb}}\,\mathbb{E}[R \mid \mathcal{F}_\tau] > c_{\text{re}}}_{\text{economics}}
\;\wedge\;
\underbrace{\Pr(R \ge h_{\min} \mid \mathcal{F}_\tau) \ge 1-\alpha}_{\text{feasibility}} ,
\end{equation}
where $\mathcal{F}_\tau$ is the information available at elapsed time
$\tau$: elapsed idle time, the \emph{current} position in the factory
calendar, and the load history of the production run that just ended.

Two aspects of \eqref{eq:rule} are deliberate. First, it is
\emph{sequential} rather than a one-shot classification at onset --- the
moment when the least information is available --- so each additional
idle minute is used as evidence. Second, the two conditions are not
redundant. Because $S$ is affine in $R$ by \eqref{eq:epsave}, the
unconstrained risk-neutral rule depends on the predictive distribution
only through its \emph{mean}, and a mean regression suffices. Constraint
\eqref{eq:c1} breaks that affinity: a shutdown whose idle period proves
shorter than $h_{\min}$ is infeasible rather than merely unprofitable,
and no expectation over $R$ expresses this. Feasibility is therefore
imposed as a chance constraint~\cite{charnes1959chance} at a stated
confidence $1-\alpha = 0.90$, using a separate classifier for
$\Pr(R\ge h_{\min})$ (AUC $0.939$, Brier $0.102$). Setting $\alpha=1$
recovers the pure risk-neutral rule and isolates the constraint's
contribution.

\paragraph{The estimand matters more than the model.}
Idle durations span four orders of magnitude, so fitting
$\log(1+R)$ is the natural choice --- and it is the wrong one here.
Back-transforming a log-scale fit estimates the conditional
\emph{median}, and Duan's smearing correction~\cite{duan1983smearing}
repairs it with a single global constant that cannot absorb a
covariate-dependent bias. At episode onset the conditional median is
$7$\,s while the conditional mean exceeds $2\,000$\,s. Since
\eqref{eq:rule} consumes the mean, the estimand must be elicited by the
loss used to fit it~\cite{gneiting2011making}; we therefore fit the raw
scale under squared error and retain the log-scale model as an ablation.

\subsection{Experimental protocol}
\label{sec:decision:protocol}

Spike rows are retained: removing $19.5\%$ of a $1$\,Hz series destroys
the correspondence between row counts and seconds on which every episode
duration depends. Episodes touching one of the record's $31$
acquisition gaps are excluded ($51$ episodes, $187.9$\,h), because their
true extent is unknown and counting them would manufacture savings out
of missing data. This leaves $3\,224$ usable episodes spanning
$474.9$\,h of idle time. The split is \emph{chronological} ($70/30$):
unlike the unsupervised comparison of Section~\ref{sec:results}, a
decision policy is deployed forward in time, and a shuffled split would
let it learn the future duty cycle. The per-day restart cap is enforced
at evaluation time, identically for every policy. All annualisation uses
the $63.4$ days the record actually covers, not its $154$-day span.

\paragraph{The status quo is the baseline.}
Of the $474.9$ usable idle hours the plant already de-energises for
$388.1$, leaving $86.8$\,h energised. Savings quoted against a
``never shut down'' hypothetical would therefore be savings the plant has
already banked; every figure below is quoted against \textsc{observed}.

\paragraph{Scenarios.}
Every coefficient is measurable except the valuation of lost production,
and $D^{\star}$ is linear in it, so asserting a value would silently
determine the conclusion. Scenarios are instead indexed by where
$D^{\star}$ falls within the observed idle-duration distribution, and the
delay valuation each implies is reported as an \emph{output}
(Table~\ref{tab:decision:scenarios}).

\begin{table}[t]
\centering
\caption{Economic scenarios, indexed by the derived break-even point.
The implied production-delay valuation is an output of
Eq.~\eqref{eq:breakeven}, not an assumption.}
\label{tab:decision:scenarios}
\begin{tabular}{lrrrr}
\hline
Scenario & $D^{\star}$ & Implied $c_{\text{dl}}$ & $c_{\text{re}}$ & Episodes above \\
         & (min) & (USD/h) & (USD) & \\
\hline
S1 low       &  10 &  1.01 & 0.075 & 228 \\
S2 moderate  &  60 &  6.09 & 0.450 &  53 \\
S3 high      & 240 & 24.34 & 1.799 &  43 \\
\hline
\end{tabular}
\end{table}

\subsection{Results}
\label{sec:decision:results}

Table~\ref{tab:decision:policies} reports annual saving against the
status quo and the share of the oracle head-room captured, on $968$
held-out episodes spanning $78.3$ days.

\begin{table}[t]
\centering
\caption{Policy comparison on held-out episodes: annual saving versus
the plant's own behaviour (USD/yr) and share of the attainable head-room
(\%). The three middle rows are ablations of the proposed policy.}
\label{tab:decision:policies}
\begin{tabular}{lrrrrrr}
\hline
& \multicolumn{2}{c}{S1 (10\,min)} & \multicolumn{2}{c}{S2 (60\,min)} & \multicolumn{2}{c}{S3 (240\,min)}\\
Policy & USD/yr & \% & USD/yr & \% & USD/yr & \% \\
\hline
Never shut down            & $-215$ & $-530$ & $-133$ & $-190$ & $+163$ & $65$ \\
Observed (status quo)      & $0$ & $0$ & $0$ & $0$ & $0$ & $0$ \\
Immediate shutdown         & $-237$ & $-584$ & $-316$ & $-451$ & $-599$ & $-241$ \\
Static break-even          & $-98$ & $-242$ & $-46$ & $-65$ & $+175$ & $70$ \\
Ski-rental (no forecast)   & $-23$ & $-56$ & $+24$ & $34$ & $+106$ & $43$ \\
\quad ablation: median estimand      & $-23$ & $-56$ & $-14$ & $-20$ & $+180$ & $72$ \\
\quad ablation: no chance constraint & $-73$ & $-180$ & $+6$ & $8$ & $+173$ & $69$ \\
\textbf{Forecast optimisation}       & $\mathbf{-3}$ & $\mathbf{-8}$ & $\mathbf{+39}$ & $\mathbf{55}$ & $\mathbf{+190}$ & $\mathbf{76}$ \\
Oracle (offline optimum)   & $+41$ & $100$ & $+70$ & $100$ & $+249$ & $100$ \\
\hline
\end{tabular}
\end{table}

\paragraph{These figures are conditional on the labelling.}
Table~\ref{tab:decision:policies} is computed on the state labels of
Section~\ref{sec:states}, in which the Viterbi path is left
unconstrained. Re-deriving the same table from labels carrying the
minimum-dwell constraint of Section~\ref{sec:flicker}---the same
machine, the same policy, the same coefficients, and a test window cut
at the same instant---changes the S2 column materially: the attainable
head-room falls from $71$\,USD/yr to $38$\,USD/yr and the proposed
policy from $+41$\,USD/yr ($57\%$ of the oracle) to
$+2.9 \pm 0.6$\,USD/yr ($7.8 \pm 1.6\%$; mean and standard deviation
over ten random seeds from Phase~V Task~12C). The physical quantities
are unmoved (idle $474.9$ vs $474.1$\,h; energised idle $86.8$ vs
$87.8$\,h; standby power $3\,747$ vs $4\,208$\,W); what moves is the
\emph{partition of idle time into episodes}, which is the unit the
decision layer bids on: enforcing a minimum dwell merges $3\,224$
episodes into $1\,312$ and raises the median episode from $12$ to
$90$\,s. Sub-ten-second ``opportunities'' created by label flicker are
not shutdown opportunities, so the smaller figure is the defensible
one, and both are reported here rather than the more favourable one
alone. The seed-to-seed standard deviation ($0.6$\,USD/yr) is small
relative to the bootstrap sampling interval ($\pm 50$\,USD/yr over
18 factory days; Phase~V Task~11B), confirming that model-fitting
variance does not explain the uncertainty.

The proposed policy captures $55$--$76\%$ of the attainable head-room in
S2 and S3 and beats the forecast-free ski-rental rule in all three
scenarios. Both ablations are material: replacing the mean estimand with
the log-scale median costs $53$\,USD/yr in S2, and removing the chance
constraint costs $33$\,USD/yr in S2 while raising minimum-off violations
from $2$ to $19$ there and from $2$ to $54$ in S1. The policy does
\emph{not} beat the status quo in S1: when the
restart cost is low the winning move is to de-energise aggressively for
long blocks, which the plant already does, and forecast noise costs more
than selectivity gains.

\subsection{Economic sensitivity}
\label{sec:decision:sensitivity}

\paragraph{Break-even is tariff-invariant absent non-energy costs.}
Sweeping the tariff over a $12\times$ range ($0.04$--$0.50$\,USD/kWh)
at $E_r = 10$\,kWh with $\ell = c_{\text{dl}} = 0$ leaves $D^{\star}$
constant at $2.669$\,h, matching $E_r/P_{\text{sb}}$ to seven
significant figures; introducing a $1.00$\,USD non-energy restart cost
makes it span $3.20$--$9.34$\,h. A tariff-versus-break-even map is
therefore uninformative for a purely electrical restart cost, and we
present the plane over restart energy and delay valuation instead, with
tariff as a third swept axis.

\begin{figure*}[t]
\centering
\includegraphics[width=\textwidth]{fig_p7_economic_plane}
\caption{The economic plane at a tariff of $0.12$\,USD/kWh, over the two
coefficients that carry information once break-even is shown to be
tariff-invariant. \textbf{(A)} The derived break-even duration
$D^{\star}$, an output of the optimisation rather than an input
threshold. \textbf{(B)} Annual saving under the proposed policy and
\textbf{(C)} energy saved over the test window, both on a diverging
scale with the zero contour drawn. \textbf{Panels B and C disagree in
sign over part of the plane}: where restarts are expensive the
profitable advice is to stop de-energising, so currency and energy move
in opposite directions and neither can be reported without the other.}
\label{fig:economic-plane}
\end{figure*}

\paragraph{Regimes.}
Figure~\ref{fig:economic-plane} maps the plane.
Over $256$ combinations of tariff, restart energy and delay valuation the
proposed policy is profitable in $249$, but only $171$ also reduce
energy: where restarts are expensive the cheapest advice is to
\emph{stop} de-energising, so currency saved and energy saved carry
opposite signs, and the two must be reported together. The regime
boundaries are
\begin{itemize}
\item the policy overtakes the status quo from $D^{\star} \approx 0.27$\,h;
\item ``never shut down'' becomes optimal from $D^{\star}\approx 27.5$\,h,
      just beyond the longest observed idle episode ($24.1$\,h);
\item the advantage over ski-rental \emph{grows} with $D^{\star}$
      ($+3.1$\,USD in S2, $+17.9$\,USD in S3 over the test window).
\end{itemize}
The last point is the practical message: when waiting is cheap a clock is
almost as good as a forecast, and forecasting earns its place precisely
where restarts are expensive.

\paragraph{Payback.}
At the measured coefficients and $D^{\star}=1$\,h the policy saves
$39$\,USD/yr per machine on the unconstrained labels of
Table~\ref{tab:decision:policies} and $2.9 \pm 0.6$\,USD/yr
(mean $\pm$ sd over ten seeds; Phase~V Task~12C) under the
minimum-dwell labelling, so a $5\,000$\,USD retrofit would pay back in
$130$ years on the first figure and not at all on the second, and a
three-year payback demands a per-machine cost below $116$\,USD in the
most favourable of the two. This follows directly from the plant already de-energising
for the long gaps: only $87$\,h of energised idle remain over $63$ days
of record, some $325$\,kWh. On pelletizer-I the framework's value is
therefore in decision quality --- it captures most of the attainable
head-room while avoiding the $20$ constraint-violating shutdowns the
status quo commits, and, in the headline scenario, its $29$ loss-making
ones --- rather than in a
large absolute saving. Section~\ref{sec:multimachine} examines whether
any other machine in the facility has a larger energised-idle fraction
and so changes this conclusion; none does.
